\documentclass[10pt,a4paper]{amsart}
\usepackage[margin=19mm]{geometry}
\usepackage{amsmath,amssymb,mathtools}
\usepackage[scr=rsfs,cal=euler]{mathalfa}
\usepackage{xfrac}
\usepackage[unicode,hidelinks,bookmarks=false]{hyperref}
\newcommand{\ie}{i.e.\ }
\newcommand{\cf}{cf.\ }
\newcommand{\resp}{respectively\ }
\newcommand{\Subsection}{\subsection}
\providecommand{\nobreakdash}{\nobreak\hbox{-}}
\setlength{\emergencystretch}{2em}
\multlinegap0pt
\usepackage{bm}
\usepackage{bbm}
\usepackage{dsfont}
\usepackage{xspace}
\usepackage{cancel}
\usepackage{mathtools}
\usepackage{amsthm}
\makeatletter
\@ifundefined{thm}{\newtheorem{thm}{Theorem}}{}
\makeatother
\makeatletter
\newcommand{\Vol}{{\rm Vol}}
\renewcommand{\epsilon}{\varepsilon}
\expandafter\ifx\csname proof\endcsname\relax
\newenvironment{proof}{\smallskip\noindent{\it Proof.}\hskip \labelsep}
\fi
\makeatother
\title[Conformal structures and rigidity of complete stable minimal]{Conformal structures and rigidity of complete stable minimal hypersurfaces}
\author{Jooyeon Park, Keomkyo Seo}
\date{Source: arXiv:2609.17078v1 (CC BY 4.0)}
\begin{document}
\maketitle

\noindent\emph{Source and licence.} Conformal structures and rigidity of complete stable minimal hypersurfaces. Version arXiv:2609.17078v1 (CC BY 4.0), distributed under the Creative Commons Attribution 4.0 International licence, as recorded in the arXiv licence metadata for this exact version and shown on the abstract page https://arxiv.org/abs/2609.17078v1. Full rights notice: licenses/arxiv-2609.17078-CC-BY-4.0.txt.\par\smallskip

\noindent\textit{Editorial scope.} Counted unit: the Thm whose source label is \texttt{thmftc2} in the article (source0.tex lines 1184--1197, source label \texttt{thmftc2}), delivered together with the article's own complete original proof of it (source0.tex lines 1199--1270, one \texttt{proof} environment of 72 lines). No mathematics is written, repaired or re-derived: statement and proof are the article's own text line for line. Required context retained uncounted: fd (lines 990--992); 9 (lines 1037--1040). Every definition, display and label of those lines is retained unchanged. Omitted from this package and not redistributed: 1--989, 993--1036, 1041--1183, 1198--1198, 1271--1690. Nothing inside a retained excerpt is omitted or altered: every retained body is delivered exactly as the article's lines are written, so no omission receipt of any kind accompanies this package. Citations: the retained excerpts carry no citation command at all, so no bibliography entry of the article is transcribed and no reference is redistributed. Reference adaptation: none was needed and none was made; every reference inside the delivered unit resolves inside it, so no unresolved reference or citation is produced. Printed slips kept literal: no printed word, symbol, hypothesis or number of the counted statement or of its proof was altered. Layout and class: the article's own class is \texttt{article} with packages including geometry, amsmath, amssymb, latexsym, amsfonts, url, graphicx, amsthm, tikz, kotex, pgfplots, color, epstopdf, hyperref; the wrapper is the lane's pinned \texttt{amsart} editorial wrapper at 10pt on A4, which restates the theorem-like environments the retained text needs and adds the article's own macro definitions that the retained text needs (\texttt{\textbackslash Vol}, \texttt{\textbackslash epsilon}), with extra wrapper packages (bm, bbm, dsfont, xspace, cancel, mathtools). The delivered numbering is the unit's own. Assets: no retained span contains a figure environment, a graphics-inclusion command, a PDF-page inclusion, a table or a TikZ picture; no image file is copied into this package, the package declares no asset dependency and no image file is redistributed with it. Licence: version arXiv:2609.17078v1 of this article is distributed under the Creative Commons Attribution 4.0 International licence, as recorded in the arXiv licence metadata for this exact version and shown on the abstract page https://arxiv.org/abs/2609.17078v1. A full rights notice is delivered at licenses/arxiv-2609.17078-CC-BY-4.0.txt. Characters: every retained excerpt is pure ASCII, so the delivered engine, XeLaTeX with its default Latin Modern text font, reports no missing character.
\begin{equation} \label{fd}
0\leq f\leq 1,\quad f\equiv 1~ \text{on}~ B_p(R),\quad f\equiv 0~\text{on}~\Sigma \backslash B_p(2R), \quad \text{and} \quad|\nabla f|\leq \frac{1}{R},
\end{equation}

\begin{align}\label{9}
\int_\Sigma &\left(2K_3-2nK_1+nK_2 +\frac{|A|^2-n(n+1)K_1+S_{\Sigma}}{2}\right) f^2|A|^2 +  \int_\Sigma |A|^2|\nabla f|^2 \nonumber\\
& \geq \frac{2}{n}\int_\Sigma f^2|\nabla |A||^2.
\end{align}

\begin{thm}\label{thmftc2}
Let $M$ be an $(n+1)$-dimensional complete Riemannian manifold whose sectional curvature $K_M$ satisfies $K_1\leq K_M\leq K_2$ for some constants $K_1$ and $K_2$, and let $\Sigma$ be a complete two-sided stable non-totally geodesic minimal hypersurface in $M$. Suppose that
$$
\lim_{R\to\infty}\frac{1}{R^2}\int_{B_p(R)}|A|^2=0,
$$
where $B_p(R)$ denotes the geodesic ball in $\Sigma$ of radius $R$ centered at some point $p\in\Sigma$. Assume further that $K_{n+1ijk}=0$ for all $1\leq i,j,k\leq n$ and 
$$
|\nabla K|^2=\sum_{i,j,k,l,m}K_{ijkl;m}^2\leq K_3^2|A|^2
$$
for some constant $K_3\geq 0$. If $\Gamma_{\Sigma}>0$, then we have
$$
\lambda_1(\Sigma)\leq \frac{n}{4}\Gamma_{\Sigma}.
$$
\end{thm}

\begin{proof}
If $\Gamma_{\Sigma}=\infty$, then the conclusion is trivial. Thus we may assume that $\Gamma_{\Sigma}<\infty$. By the variational definition of $\lambda_1(\Sigma)$, for every nonzero compactly supported Lipschitz function $\phi$ on $\Sigma$,
    \begin{equation}\label{10}
        \lambda_1(\Sigma) \leq \frac{\int_\Sigma|\nabla \phi|^2}{\int_\Sigma \phi^2}.
    \end{equation}
Choose $\phi=|A|f$ in \eqref{10}, where $f$ is the cut-off function defined in \eqref{fd}. Then
  \begin{align}\label{11}
\lambda_1(\Sigma)\int_\Sigma f^2|A|^2
&\le \int_\Sigma |\nabla (|A|f)|^2 \nonumber\\
&= \int_\Sigma f^2|\nabla |A||^2
 + \int_\Sigma |A|^2|\nabla f|^2
 + 2\int_\Sigma f|A|\langle \nabla f,\nabla |A| \rangle.
\end{align}
Moreover, Young's inequality gives
  \begin{equation}\label{12}
        2\int_\Sigma f|A|\langle \nabla f, \nabla |A|\rangle
        \leq \epsilon\int_\Sigma f^2|\nabla |A||^2 + \frac{1}{\epsilon} \int_\Sigma |A|^2|\nabla f|^2
    \end{equation}
    for every $\epsilon>0$. Plugging (\ref{12}) into (\ref{11}), we obtain
  \begin{equation}\label{13}
\lambda_1(\Sigma)\int_\Sigma f^2|A|^2
\le
\left(1+\epsilon\right)\int_\Sigma f^2|\nabla |A||^2
+ \left(1+\frac{1}{\epsilon}\right)\int_\Sigma |A|^2|\nabla f|^2.
\end{equation}
On the other hand, using the definition of $\Gamma_\Sigma$ and  (\ref{9}) gives
      \begin{align}\label{14}
        \frac{\Gamma_{\Sigma}}{2}&\int_\Sigma f^2|A|^2 + \int_\Sigma |A|^2|\nabla f|^2 \nonumber\\
        &\geq \frac{1}{2}\int_\Sigma(-(n^2+5n)K_1 + 2nK_2 + 4K_3 + |A|^2 + S_{\Sigma}) f^2|A|^2 + \int_\Sigma |A|^2|\nabla f|^2 \nonumber \\
        & = \int_\Sigma\left(2K_3-2nK_1+nK_2+\frac{|A|^2-n(n+1)K_1+S_{\Sigma}}{2}\right) f^2|A|^2 + \int_\Sigma |A|^2|\nabla f|^2 \nonumber\\
        & \geq \frac{2}{n}\int_\Sigma f^2|\nabla |A||^2.
\end{align}
If $\lambda_1(\Sigma)=0$, then the desired estimate holds trivially since $\Gamma_\Sigma>0$. Hence we may assume that $\lambda_1(\Sigma)>0$. Combining \eqref{13} and \eqref{14} and using the assumption that $\Gamma_\Sigma>0$, we see that
  \begin{equation*}
\left\{
1+\frac{\Gamma_\Sigma}{2\lambda_1(\Sigma)}\left(1+\frac{1}{\epsilon}\right)
\right\}
\int_\Sigma |A|^2|\nabla f|^2
\ge
\left\{
\frac{2}{n}-\frac{\Gamma_\Sigma}{2\lambda_1(\Sigma)}(1+\epsilon)
\right\}
\int_\Sigma f^2|\nabla |A||^2.
\end{equation*}
Now suppose, for contradiction, that
$$
\lambda_1(\Sigma)>\frac{n}{4}\Gamma_\Sigma.
$$
Then we may choose $\epsilon>0$ sufficiently small so that
$$
\frac{2}{n}-\frac{\Gamma_\Sigma}{2\lambda_1(\Sigma)}(1+\epsilon)>0.
$$
Since $|\nabla f|\le 1/R$ and $\operatorname{supp}f\subset B_p(2R)$, the growth assumption gives
$$
\int_\Sigma |A|^2|\nabla f|^2
\le
\frac{1}{R^2}\int_{B_p(2R)}|A|^2\longrightarrow0.
$$
  Because $f\equiv1$ on $B_p(R)$ and the coefficient on the right-hand side is strictly positive, we obtain $\nabla|A|\equiv0$. Hence $|A|\equiv c$ for some constant $c\ge0$.

If $c>0$, applying \eqref{10} to the cut-off function $f$ itself gives
$$
\lambda_1(\Sigma)\int_\Sigma f^2
\le
\int_\Sigma|\nabla f|^2
\le
\frac{1}{R^2}\Vol(B_p(2R))
=
\frac{1}{c^2R^2}\int_{B_p(2R)}|A|^2\longrightarrow0.
$$
Since $f\equiv1$ on $B_p(R)$, the left-hand side is bounded below by $\lambda_1(\Sigma)\Vol(B_p(R_0))>0$ for any fixed $R_0$ and all sufficiently large $R$, a contradiction. Thus $c=0$, contrary to the assumption that $\Sigma$ is not totally geodesic. This completes the proof.
\end{proof}

\end{document}
